% =====================================================================
%  E14: E13's main sets re-run with all 35 rule templates (Yang et al.'s
%  five and the 30 new E4b templates): means from every template, and the
%  same rules re-applied with every score computed on all 35 templates.
%  Tables and figure are generated by scripts/e14_report.py
%  (paper/tables/e14_*.tex, paper/figures/e14_groups.pdf).
% =====================================================================
\documentclass[10pt,a4paper]{article}

\usepackage[utf8]{inputenc}
\usepackage[T1]{fontenc}
\usepackage[margin=1.8cm]{geometry}
\usepackage{amsmath,amssymb}
\usepackage{booktabs,tabularx,array}
\usepackage{enumitem}
\usepackage{xcolor}
\usepackage{caption}
\usepackage{float}
\usepackage{graphicx}
\usepackage[colorlinks=true,linkcolor=blue!50!black,urlcolor=blue!60!black]{hyperref}

\setlist{nosep,leftmargin=*}
\setlength{\parindent}{0pt}
\setlength{\parskip}{5pt}
\renewcommand{\arraystretch}{1.12}
\captionsetup{font=small,labelfont=bf}

\title{E14: The Main Head Sets with All 35 Wordings\\\large Lexical task heads, simple tasks}
\author{Filippo Merlo}
\date{12 November 2026}

\begin{document}
\maketitle
\vspace{-1.5em}

%INSHORT

% ---------------------------------------------------------------------
\section{Why}

Every set of E13 was chosen, and every patch was made, with Yang et al.'s five rule templates only: the word lists, D, B, T and the generic score came from their correct prompts, Davidson et al.'s causal score from the four of them that have an instruction before ``Q: \ldots\ A:'', and each set was patched with the five templates' means. The 30 new templates per task written for E4b (proposal-simple, Sec.~16) were not used. E14 re-runs E13's main conditions with all 35 templates, in two ways: the same sets with the means of all 35 templates, and the same rules with every score recomputed on all 35.

% ---------------------------------------------------------------------
\section{Design}

\begin{enumerate}
\item \textbf{Templates}: per task Yang et al.'s five and the 30 new ones, in six groups of five: form (cloze, question--answer, label) $\times$ whether the template names the task. The new templates were generated by GPT-5.5 and checked by script only; the planned hand check was skipped (the author's decision), so all 30 per task are used as generated.
\item \textbf{E0 on the new templates}: greedy replies, flexible score, as for Yang's templates, on the first 500 items (in file order) of each half of E4's split. Every later step takes at most the first 100 correct prompts of a half, so the cap changes nothing where a template has 100 correct among them.
\item \textbf{Conditions} (E13's best and the references, at 40 and 80 heads): DBT, Causal and Causal $-$ generic, one set per task and one set for all tasks; random sets with as many heads in every layer as the DBT-80 set chosen on all 35 (five draws). Left out: BT (within 0.02 of DBT in E13), the other filter combinations, Yang terms, DBT $-$ generic, Davidson's summed vector, and 20 and 60 heads.
\item \textbf{Heads chosen on Yang's 5}: E13's sets, unchanged.
\item \textbf{Heads chosen on all 35}: E13's rules, every score on the choosing half of the 35 templates. Word lists and D: the J-head top 10 at the last token of every template's correct prompts (at most 100 each, so up to 3{,}500 per task instead of 500), pooled. B and T: a new attention pass over the same prompts, T on the positions of list words outside the query. Causal: Davidson et al.'s score on the 14 templates of the question--answer form (Yang's four and the 10 new ones), 25 mean and 5 scoring prompts each; the cloze and label templates have no instruction to replace. The other tasks' new question--answer instructions join the pool of the ``other task'' uninformative kind. Generic heads: the J-head generic score on the 35 templates, cut among the fold's new Davidson top 30 by E13's two-group split.
\item \textbf{Means}: every set is patched with each template's own mean, one at a time (the first 100 correct prompts of the template on the choosing half; a template with no correct prompt there gives none), and the control task's mean as in E13. Check~1 only, $\alpha\in\{1,2\}$, replies of at most 16 tokens, both folds. Task accuracy = mean over the templates' means; $\alpha$ is cross-fitted on it (the better on the other fold, ties to 1). The breakdown by template group is read at the same $\alpha$. Random draws: task means only, at the $\alpha$ of their DBT set.
\item \textbf{Check}: the sets chosen on Yang's 5, with Yang's five means only, must give E13's numbers.
\end{enumerate}

% ---------------------------------------------------------------------
\section{Templates and sets}

\begin{table}[H]\centering\scriptsize
\resizebox{\linewidth}{!}{% Generated by scripts/e14_report.py.
\begin{tabular}{@{}lccccccccc@{}}
\toprule
 & \multicolumn{7}{c}{E0 accuracy (flexible), mean of the group's templates} & \multicolumn{2}{c}{Templates} \\
\cmidrule(lr){2-8}\cmidrule(lr){9-10}
Task & \begin{tabular}[b]{@{}c@{}}Yang's 5\end{tabular} & \begin{tabular}[b]{@{}c@{}}cloze,\\named\end{tabular} & \begin{tabular}[b]{@{}c@{}}cloze,\\not named\end{tabular} & \begin{tabular}[b]{@{}c@{}}Q--A,\\named\end{tabular} & \begin{tabular}[b]{@{}c@{}}Q--A,\\not named\end{tabular} & \begin{tabular}[b]{@{}c@{}}label,\\named\end{tabular} & \begin{tabular}[b]{@{}c@{}}label,\\not named\end{tabular} & $<30$ & 0 \\
\midrule
antonym & 0.68 & 0.37 & 0.52 & 0.67 & 0.68 & 0.67 & 0.67 & 1 & 0 \\
synonym & 0.45 & 0.30 & 0.30 & 0.48 & 0.49 & 0.20 & 0.23 & 2 & 1 \\
singular--plural & 0.98 & 0.58 & 0.23 & 0.99 & 0.96 & 0.99 & 0.94 & 5 & 3 \\
present--past & 0.96 & 0.39 & 0.39 & 0.98 & 0.93 & 0.99 & 0.65 & 6 & 4 \\
country--capital & 0.97 & 0.70 & 0.77 & 0.97 & 0.94 & 0.97 & 0.93 & 2 & 1 \\
country--currency & 0.72 & 0.45 & 0.66 & 0.68 & 0.45 & 0.66 & 0.10 & 9 & 5 \\
landmark--country & 0.85 & 0.55 & 0.66 & 0.83 & 0.82 & 0.70 & 0.50 & 1 & 1 \\
park--country & 0.79 & 0.57 & 0.32 & 0.89 & 0.84 & 0.89 & 0.61 & 1 & 0 \\
product--company & 0.73 & 0.58 & 0.57 & 0.73 & 0.68 & 0.74 & 0.69 & 0 & 0 \\
next-item & 0.42 & 0.44 & 0.32 & 0.39 & 0.20 & 0.50 & 0.64 & 12 & 1 \\
prev-item & 0.53 & 0.29 & 0.32 & 0.54 & 0.35 & 0.46 & 0.42 & 11 & 0 \\
English--French & 0.73 & 0.53 & 0.22 & 0.78 & 0.64 & 0.73 & 0.71 & 1 & 0 \\
English--German & 0.67 & 0.49 & 0.22 & 0.68 & 0.46 & 0.73 & 0.40 & 5 & 3 \\
English--Spanish & 0.79 & 0.59 & 0.46 & 0.82 & 0.74 & 0.71 & 0.82 & 1 & 1 \\
\midrule
\textit{Mean} & 0.73 & 0.49 & 0.43 & 0.75 & 0.65 & 0.71 & 0.59 & 57 & 20 \\
\bottomrule
\end{tabular}
}
\caption{E0 on the 35 templates: flexible accuracy without a patch (zero-shot), mean of each group's five templates, on the first 500 items of each half (the E0 items of the new templates). $<30$: templates with fewer than 30 correct prompts on at least one half, so a mean from few prompts on that fold; 0: templates with none on at least one half, which give no mean there.}
\label{tab:e0}
\end{table}

The question--answer and label templates that name the task do about as well as Yang et al.'s five. The cloze templates, and the templates that do not name the task, do worse. Where a cloze template fails, the model goes on with the sentence before the answer or instead of it: ``Converting adapt to past tense gives \emph{us} adapted'', ``The form of wallet for two or more is \emph{called} a wallet'', ``\ldots\ the English word seminary is \emph{translated} to German as \ldots''. The flexible score needs the reply to begin with the answer, so such replies count as wrong even when the answer comes later. A few label templates fail outright (``State: Afghanistan / Primary tender:'' gets ``1'' or a year). The templates with no correct prompt on a half are nearly all cloze templates or label templates that do not name the task, and those with few correct prompts are mostly on the six small tasks (the two morphology tasks, country--capital, country--currency, next and previous item), which have 100--150 items per half.

The word lists from the 35 templates keep 19--27 of the 30 words of E13's lists, from five to seven times as many prompts.

\begin{table}[H]\centering\small
% Generated by scripts/e14_report.py.
\begin{tabular}{@{}lcccccc@{}}
\toprule
 & \multicolumn{3}{c}{One set per task} & \multicolumn{3}{c}{One set for all 14 tasks} \\
\cmidrule(lr){2-4}\cmidrule(lr){5-7}
Heads (N) & DBT & Causal & Causal $-$ generic & DBT & Causal & Causal $-$ generic \\
\midrule
40 & 29.3 & 23.6 & 22.0 & 35.5 & 33.0 & 30.0 \\
80 & 58.7 & 46.6 & 44.6 & 61.0 & 65.5 & 63.0 \\
100 & 73.1 & 58.5 & 55.7 & 79.0 & 83.5 & 77.0 \\
120 & 88.1 & 68.9 & 66.3 & 93.0 & 94.5 & 91.0 \\
\bottomrule
\end{tabular}

\caption{Heads of each set chosen on all 35 templates that are also in the same set chosen on Yang's 5 (E13), mean over tasks and folds.}
\label{tab:overlap}
\end{table}

Re-choosing on all 35 templates changes DBT least: on every task about three quarters of the heads stay, at both sizes. Causal changes more, and unevenly. On the three translation tasks and on country--currency, at 40 heads, half or fewer of the heads stay; on antonym, singular--plural and the fact tasks most of them do. The shared sets change little: the rank averaging over 14 tasks evens out the per-task changes. Davidson et al.'s top 30 over all tasks keeps 26 (fold~1) and 25 (fold~2) of E13's heads. The generic heads cut from it are much the same: L13H31, L14H29, L15H5, L16H29 and L31H14 in both folds and L28H10 in fold~1, as in E13. L12H9, L14H11 and L15H7 are no longer marked, and L13H6 is newly marked in both folds (with L13H25 in fold~1 and L29H23 in fold~2).

%RESULTS

\end{document}
